\needspace{0.25\textheight}
\section{Критический разбор: границы переноса на полное уравнение}\label{sec:critique}

\subsection{Структурный словарь}

Сведём соответствие между моделью \eqref{eq:epd} и полным уравнением
\eqref{eq:ns-vort} (см.~\cite{arnoldkhesin,marsdenratiu}):

\begin{table}[htbp]
\centering
\small
\caption{Структурный словарь «модель $\leftrightarrow$ полное уравнение».}
\label{tab:dictionary}
\begin{tabularx}{\textwidth}{@{}l X l@{}}
\toprule
Модель & Содержание & Полное уравнение \\
\midrule
$\Om \in \so^{*}$ & амплитудный вектор вихревой трубки &
  $\bwn(x)$, бесконечномерно \\
$\Om\times\bbI^{-1}\Om$ & коалгебраический транспорт (коммутатор) &
  $-[u,\bwn] = \nabla\times(u\times\bwn)$ \\
$-\nu\bbL\Om$, $\bbL\succeq0$ & диссипация через положительный символ &
  $\nu\Delta\bwn$ \\
$C = \|\Om\|^2$ & Казимир & $\tfrac12\|\bwn\|_{L^2}^2$ (энстрофия) \\
$\mathcal H$ & энергия & $\tfrac12\|u\|_{L^2}^2$ \\
$S^2_{\|\Om_0\|} \subset \so^{*}$ & коалгебраическая орбита &
  орбиты группы диффеоморфизмов \\
\bottomrule
\end{tabularx}
\end{table}

\subsection{Почему модель работает: сферичность орбит}

\begin{proposition}[орбиты $\so^{*}$]\label{prop:orbits}
Коалгебраические орбиты действия $\mathrm{SO}(3)$ на $\so^{*}
\cong \R^3$~--- в точности сферы $\{\|\Om\| = \mathrm{const}\}$.
Следовательно любая радиальная функция, в частности $C = \|\Om\|^2$,~---
Казимир скобки \eqref{eq:lp}, и сохранение $\|\Om\|$ при $\nu=0$
является \emph{чисто геометрическим} фактом.
\end{proposition}

\begin{proof}
Инфинитезимальное коалгебраическое действие вычислено в
\eqref{eq:adstar}: $\adst_{\xi}\Om = \Om\times\xi$. Касательное
пространство орбиты в точке $\Om$ есть
$\{\Om\times\xi : \xi \in \R^3\} = (\operatorname{span}\{\Om\})^{\perp}$~--- касательная плоскость к сфере радиуса $\|\Om\|$ в точке $\Om$.
Значит орбита~--- открытое подмногообразие соответствующей сферы; будучи
орбитой действия компактной группы, она компактна и потому совпадает
со всей сферой (сфера связна). Радиальные функции постоянны на орбитах,
т.е. коммутируют со скобкой \eqref{eq:lp} в смысле
предложения~\ref{prop:poisson}.
\end{proof}

В этом~--- точное структурное объяснение всей конструкции: сферичность
орбит \emph{гарантирует} закон Казимира. Для двумерной идеальной
жидкости аналогичное утверждение истинно: энстрофия
$\tfrac12\|\bwn\|_{L^2}^2$~--- подлинный Казимир бесконечномерной
скобки Ли--Пуансона группы сохраняющих площадь диффеоморфизмов, и
теорема Ладыженской \eqref{eq:enstrophy} есть её вязкий аналог
\cite{arnoldkhesin,ladyzh}. Для трёхмерной идеальной жидкости Казимиром
служит только гелисичность, а не энстрофия: орбиты несферичны, и член
растяжения $(\bwn\cdot\nabla)\bu = S\bwn$ (лемма~\ref{lem:strain})
выводит норму $\|\bwn\|_{L^2}$ с орбитально-инвариантных множеств.
Баланс энстрофии полного уравнения
\begin{equation}\label{eq:full-enstrophy}
  \frac{d}{dt}\,\frac12\|\bwn\|_{L^2}^2
  = \int \langle S\bwn, \bwn\rangle\,dx \;-\; \nu\|\nabla\bwn\|_{L^2}^2
\end{equation}
содержит штрейн-член $\int\langle S\bwn,\bwn\rangle dx$, который модель
обнуляет \emph{по построению}: в \eqref{eq:epd} ему просто негде жить~---
кососимметричная нелинейность $\Om\times\bbI^{-1}\Om$ структурно
не способна его породить (теорема~\ref{th:main}).

\subsection{Условная теорема о тензоре деформации}

Сделаем границу метода количественной. Расширим модель явным
штрейн-членом:
\begin{equation}\label{eq:epd-strain}
  \dot\Om = \Om\times\bbI^{-1}\Om + S(t)\,\Om - \nu\,\bbL\,\Om,
\end{equation}
где $S(t) = S(t)^{\mathsf T}$~--- произвольные ограниченные измеримые
симметричные матрицы (экзогенный тензор деформации).

\begin{theorem}[условная]\label{th:strain}
Вдоль решений \eqref{eq:epd-strain} при \textup{(H1)} и
$\lambda_{\min} := \lambda_{\min}(\bbL) > 0$
\begin{equation}\label{eq:strain-bound}
  \|\Om(t)\| \;\le\; \|\Om_0\|\,
  \exp\!\Bigl(\int_0^t \|S(s)\|\,ds \;-\; \nu\lambda_{\min}\,t\Bigr)
  \qquad \forall\, t \ge 0 .
\end{equation}
В частности, если штрейн интегрируем с запасом:
$\int_0^t\|S(s)\|ds \le \nu\lambda_{\min}t + K$ на $[0,T]$, то
$\|\Om(t)\| \le e^{K}\|\Om_0\|$~--- сингулярность невозможна.
\end{theorem}

\begin{proof}
Дифференцируем $C = \|\Om\|^2$ вдоль \eqref{eq:epd-strain}:
\[
  \dot C = 2\underbrace{\langle\Om,\Om\times\bbI^{-1}\Om\rangle}_{=0}
  + 2\langle S(t)\Om,\Om\rangle - 2\nu\langle\Om,\bbL\Om\rangle
  \le \bigl(2\|S(t)\| - 2\nu\lambda_{\min}\bigr)\,C ,
\]
где обнуление перекрёстного члена~--- как в теореме~\ref{th:casimir},
а неравенство~--- Коши--Буняковский с $\|S\| := \|S\|_{2}$ и оценка
Рэлея для $\bbL$. Дифференциальное неравенство
$\dot C \le \beta(t)C$ интегрируется (Грёнуолл):
$C(t) \le C(0)\exp\int_0^t\beta(s)ds$, что после извлечения корня даёт
\eqref{eq:strain-bound}. Следствие о $e^K\|\Om_0\|$~--- подстановка
условия интегрируемости.
\end{proof}

\begin{corollary}[аналогия с BKM/CKN]\label{cor:bkm}
Структура теоремы~\ref{th:strain}~--- точный аналог классических
критериев регулярности полного уравнения: критерий
Бил--Като--Маджды связывает отсутствие сингулярности на $[0,T]$ с
конечностью $\int_0^T\|\bwn\|_{L^\infty}dt$ \cite{bkm}, критерий
Каффарелли--Кона--Ниренберга~--- с малостью интеграла от
характеристики $A(t) \sim \bigl(1 + \|\nabla\bu\|_{L^\infty}\bigr)$
\cite{ckn}. В модели роль $\|\bwn\|_{L^\infty}$ играет $\|S(t)\|$:
\emph{регулярность определяется интегрируемой силой растяжения}.
Численно оценка \eqref{eq:strain-bound} точна: при изотропном
$S \equiv s_0\,\mathrm{Id}$ выполняется $\|\Om(t)\| =
\|\Om_0\|e^{(s_0-\nu\lambda_{\min})t}$ с \emph{равенством} в
\eqref{eq:strain-bound} (тест S1b в \S\ref{sec:numerics}), а при
$s_0 > \nu\lambda_{\min}$ реализуется экспоненциальный рост~--- модель
с несжимаемым штрейном не имеет априорной оценки, что и есть
честная граница метода.
\end{corollary}

\subsection{Что доказано и что остаётся открытым}

\begin{enumerate}[label=\textup{(\arabic*)}, leftmargin=2.2em, itemsep=2pt]
  \item \emph{Доказано строго} (внутри модели): все результаты
        \S\ref{sec:geom}--\S\ref{sec:noblowup}~--- геометрия
        $\RP^3\cong\mathrm{SO}(3)$, вывод \eqref{eq:epd} из
        вариационного принципа, скобка Ли--Пуансона с тождеством
        Якоби, инварианты и оценки, периодичность, явные решения,
        отсутствие сингулярностей \eqref{eq:noblowup}, оценка МакВини
        с оптимальной константой. Всё это подтверждено численно
        (23/23, \S\ref{sec:numerics}).
  \item \emph{Переносится на полное уравнение}: только каскад
        $2D \to 1D \to 0D$ в той части, где он опирается на
        действительные теоремы (Ладыженская, Коул--Хопф); структурная
        аналогия Казимира верна в $2D$, где энстрофия~--- подлинный
        Казимир \cite{arnoldkhesin}.
  \item \emph{Не переносится}: заключение \eqref{eq:noblowup} о
        трёхмерном уравнении. Причина~--- замыкание петли: в полном
        уравнении $S(t) = \operatorname{sym}(\nabla\bu)$ не экзогенен,
        а определяется тем же полем $\bu$ (через уравнение для
        $\bwn$ и оператор Био--Савара). Модель размыкает эту петлю,
        постулируя $S$ заданным, и потому её теорема~\ref{th:strain}
       ~--- условное утверждение вида «\emph{если} штрейн
        интегрируем, \emph{то} сингулярности нет». Именно
        неэкзогенность $S$ делает проблему тысячелетия открытой
        \cite{fefferman,ckn,bkm}; сведение её к конечной размерности
        само по себе не даёт контроля $\|S(t)\|_{L^\infty}$.
  \item \emph{Технические оговорки}: компактность $\RP^3$
        ограничивает состояние, но не скорости (в модели скорости
        ограничены априорной оценкой, в полном уравнении~--- нет);
        энергетическая монотонность требует \textup{(H2)}
        (замечание~\ref{rem:commutator}); $\bbP(\Om)$ не определена в
        нуле, куда траектория сходится при $\bbL \succ 0$.
\end{enumerate}

Итог: модель $P^3 \to \so$~--- строго верифицируемый конечномерный
каркас геометрии турбулентной завихренности, точный в $2D$-пределе и
условно точный в $3D$; его главный вывод для полного уравнения~---
не «сингулярности нет», а «сингулярность эквивалентна
неинтегрируемости штрейна», что совпадает с направлением классических
критериев \cite{bkm,ckn}.
